\documentclass[ignorenonframetext,]{beamer}
\setbeamertemplate{caption}[numbered]
\setbeamertemplate{caption label separator}{: }
\setbeamercolor{caption name}{fg=normal text.fg}
\beamertemplatenavigationsymbolsempty
\usepackage{lmodern}
\usepackage{amssymb,amsmath}
\usepackage{ifxetex,ifluatex}
\usepackage{bm}
\usepackage{extarrows}
\usepackage{fixltx2e} % provides \textsubscript
\ifnum 0\ifxetex 1\fi\ifluatex 1\fi=0 % if pdftex
\usepackage[T1]{fontenc}
\usepackage[utf8]{inputenc}
\else % if luatex or xelatex
\ifxetex
\usepackage{mathspec}
\else
\usepackage{fontspec}
\fi
\defaultfontfeatures{Ligatures=TeX,Scale=MatchLowercase}
\fi
\usetheme{madrid}
\usecolortheme{orchid}
% use upquote if available, for straight quotes in verbatim environments
\IfFileExists{upquote.sty}{\usepackage{upquote}}{}
% use microtype if available
\IfFileExists{microtype.sty}{%
\usepackage{microtype}
\UseMicrotypeSet[protrusion]{basicmath} % disable protrusion for tt fonts
}{}
\newif\ifbibliography
\usepackage{natbib}
\usepackage{graphicx,grffile}
\makeatletter
\def\maxwidth{\ifdim\Gin@nat@width>\linewidth\linewidth\else\Gin@nat@width\fi}
\def\maxheight{\ifdim\Gin@nat@height>\textheight0.8\textheight\else\Gin@nat@height\fi}
\makeatother
% Scale images if necessary, so that they will not overflow the page
% margins by default, and it is still possible to overwrite the defaults
% using explicit options in \includegraphics[width, height, ...]{}
\setkeys{Gin}{width=\maxwidth,height=\maxheight,keepaspectratio}
\definecolor{azure}{rgb}{0, 0.4, 0.8}
\definecolor{crimson}{rgb}{0.8, 0.4, 0}
\definecolor{gold}{rgb}{0.72, 0.53, 0.04}
\definecolor{ciel}{rgb}{0.39, 0.58, 0.93}
\definecolor{dgreen}{rgb}{0.01, 0.75, 0.24}
% Prevent slide breaks in the middle of a paragraph:
\widowpenalties 1 10000
\raggedbottom

\setlength{\parindent}{0pt}
\setlength{\parskip}{6pt plus 2pt minus 1pt}
\setlength{\emergencystretch}{3em}  % prevent overfull lines
\providecommand{\tightlist}{%
\setlength{\itemsep}{0pt}\setlength{\parskip}{0pt}}
\setcounter{secnumdepth}{0}
\usepackage{bbm}
\usepackage{hhline}
\newcommand{\E}[1]{{\mathbb E}\left[#1\right]}
\newcommand{\Var}[1]{{\mathrm Var}\left(#1\right)}
\usepackage{beamerthemesplit}
\usepackage{amsmath}
\usepackage{amsthm}
\usepackage{kotex}
\usepackage{booktabs}
\usepackage{multirow}
\usepackage{makecell}
\usepackage{array}
\usepackage{arydshln}
\usepackage{caption}
\usepackage{comment}
\newcommand{\parallelsum}{\mathbin{\!/\mkern-5mu/\!}}
\newcommand\tab[1][1cm]{\hspace*{#1}}
\captionsetup[table]{belowskip=0pt,aboveskip=0pt}
\captionsetup[figure]{belowskip=0pt,aboveskip=0pt}
\newtheorem{model}{Model}

\title{Multiple Linear Regression}

\author[Himchan Jeong]{\large{Himchan Jeong, PhD, FSA}}

%\institute{Simon Fraser University}

\date{ }

\begin{document}
\frame{\titlepage}



\begin{frame}{Motivation}

Recall the basic assumptions of the linear models where $y_1,\ldots,y_m$ are the observed response variable and $\mathbf{x}_1,\ldots, \mathbf{x}_m \in \mathbb{R}^{p \times 1}$ are the individual covariates  (here $Z \in \mathbb{R}^{a \times b}$ means that $Z$ is a $(a \times b)$ real-valued matrix):
\begin{enumerate}
	\item $E[y_i] = \mathbf{x}_i' \bm{\beta}$ for $\bm{\beta} = (\beta_1, \ldots, \beta_p) \in \mathbb{R}^{p\times 1}$,
	\item $ \{\mathbf{x}_1, \mathbf{x}_2, \ldots, \mathbf{x}_m \}$ are non-random,
	\item $Var(y_i) = \sigma^2$,
	\item $\{ y_1, y_2, \ldots, y_m  \}$ are independent,
	\item $\{ y_1, y_2, \ldots, y_m \}$ are normally distributed.
\end{enumerate}
Such situations can be written in short as follows:
\begin{equation} \label{eq:linmod}
y_i = \mathbf{x}_i' \bm{\beta}  + \epsilon_i, \quad \epsilon_i \overset{i.i.d.}\sim \mathcal{N}(0, \sigma^2).
\end{equation}

Note that the first components of $ \{\mathbf{x}_1, \mathbf{x}_2, \ldots, \mathbf{x}_m \}$ are all 1 so that $\beta_1$ corresponds to the \emph{intercept} term in the linear models.

\end{frame}

\begin{frame}{Model assumptions and notation}

Here is a general form of the linear models with matrix notation:

\begin{equation} \label{eq:linmodgen}
\begin{aligned}
&\bm{y} = \bm{X\beta + \epsilon}, \ \bm{\epsilon} \sim \mathcal{N}(\bm{0, \Sigma}=\sigma^2 I_m),
\end{aligned}
\end{equation}
where $\bm{X} = \begin{pmatrix} \mathbf{x}'_1 \\ \mathbf{x}'_2 \\ \vdots \\ \mathbf{x}_m' \end{pmatrix}\in \mathbb{R}^{m \times p}$ is the covariate matrix for $m$ observations, $\bm{\Sigma} \in \mathbb{R}^{m \times m}$ is the covariance matrix of the white noise $\bm{\epsilon} \in \mathbb{R}^{m \times 1}$. 

The marginal mean and variance of $\mathbf{y} = (y_1, y_2, \ldots, y_m) \in \mathbb{R}^{m \times 1}$ are given as follows:
\begin{itemize}
	\item $E[\bm{y}] = \bm{X\beta} \in \mathbb{R}^{m \times 1} $, $Var(\bm{y}) = \sigma^2 I_m \in \mathbb{R}^{m \times m}$.
\end{itemize}



\end{frame}

\begin{frame}[allowframebreaks]
\frametitle{Point Estimation}
\begin{theorem}[Estimation of $\bm{\beta}$ - OLS]
If $X  \in \mathbb{R}^{m \times p}$ is of full rank, then the ordinary least squares (OLS) estimator of $\bm{\beta}$ is given as follows:
$$
\hat{\bm{\beta}}^{OLS}=\bm{(X'X)^{-1} X'y}.
$$
\end{theorem}

\begin{proof}
See \url{https://en.wikipedia.org/wiki/Proofs_involving_ordinary_least_squares}.
\end{proof}

\framebreak
\begin{theorem}[Estimation of $\bm{\beta}$ - MLE]
If $X \in \mathbb{R}^{m \times p}$ is of full rank and assume that $\bm{\epsilon} \sim \mathcal{N}(\bm{0, \Sigma}=\sigma^2 I_m)$, then the maximum likelihood estimator (MLE) of $\bm{\beta}=(\beta_1, \ldots, \beta_p) \in \mathbb{R}^{p \times 1}$ is given as follows:
$$
\hat{\bm{\beta}}=\bm{(X'X)^{-1} X'y}.
$$
\end{theorem}

\begin{proof}
$$
\begin{aligned}
\ell(\beta) &= \sum_{i=1}^m \log f_{y}(y_i;\mathbf{x}_i, \beta) =   -\sum_{i=1}^m  \left( \frac{(y_{i}-  \mathbf{x}'_i\beta)^2}{\sigma^2} \right)  + \text{const}.
\end{aligned}
$$
Remaining steps are to be explored together.
\end{proof}

\framebreak

\begin{theorem}[Estimation of $\sigma^2$]
If $X \in \mathbb{R}^{m \times p}$ is of full rank, then the unbiased estimator of $\sigma^2$ is given as follows:
$$
\hat{\sigma}^2= \frac{1}{m-p} y' M y,
$$
where $M = I_m - X(X'X)^{-1}X'$ and  $p$ is the rank of $X$ (intuitively speaking, the number of independent covariates in $X$).
\end{theorem}
Note that $E[\hat{\sigma}^2] = \frac{m-p}{m-p} \sigma^2=  \sigma^2$.

\begin{proof}
$$
\begin{aligned}
E[y' M y] &\overset{(a)}{=} E[\bm{\epsilon}' M \bm{\epsilon}] \overset{(b)}{=}  E[\text{tr}(\bm{\epsilon}' M \bm{\epsilon})]\overset{(c)}{=}  E[\text{tr}(M\bm{\epsilon}  \bm{\epsilon}')]\\ 
&\overset{(d)}{=}  \text{tr}(E[M\bm{\epsilon}  \bm{\epsilon}']) 
\overset{(e)}{=} \text{tr}(M \cdot E[\bm{\epsilon}  \bm{\epsilon}'])\overset{(f)}{=} \text{tr}(M \cdot \sigma^2 I_m) \\
&\overset{(g)}{=} \sigma^2 \text{tr}(M) \overset{(h)}{=} \sigma^2 (m-p).
\end{aligned}
$$
\end{proof}

\framebreak
Here $\text{Tr}(A)$ (trace) means the diagonal sum of a matrix $A$.
\begin{itemize}
\item (a): By definition of $M$.
\item (b): $\bm{\epsilon}' M \bm{\epsilon}$ is a $(1\times1)$ matrix.
\item (c): By definition of trace (\url{https://en.wikipedia.org/wiki/Trace_(linear_algebra)}).
\item (d): By definition of trace (\url{https://math.stackexchange.com/questions/2228398/trace-trick-for-expectations-of-quadratic-forms}).
\item (e): $M$ and $\bm{\epsilon}$ are independent.
\item (f): By the assumption in  \eqref{eq:linmodgen}. Note that $E[\bm{\epsilon}\bm{\epsilon}']=Var(\bm{\epsilon})$ (why?)
\item (g): Obvious.
\item (h): $\text{tr}(M) = \text{tr}(I_m) - \text{tr}(X(X'X)^{-1}X') =  \text{tr}(I_m) - \text{tr}(X'X(X'X)^{-1}) = m - p$ (why?)
\end{itemize}


\end{frame}


\begin{frame}{Confidence about regression models}

Suppose we have two regression models from different data, where the fitted regression line is $\hat{y}_i =  1 + 0.5x_i$ for both.

     \centering
     \includegraphics[trim={0cm 0.5cm 0cm 0.8cm},clip,height=8cm, width=12cm]{regplots.pdf}

Shall we have the same level of `confidence' on both regression models?

\end{frame}

%\begin{theorem}[Prediction of $\bm{y}$ with known $\bm{V}$ and updated $\bm{X}_{i,*}$]
%Suppose that for the subject $i$, we observed a new covariate vector $\bm{X}_{i,*}$. In this case, the predicted value of $\bm{y}_{i,*}$ is given as follows:
%$$
%\begin{aligned}
%\hat{\bm{y}}_{i,*}:&= \bm{X}_{i,*}\hat{\bm{\beta}} + \bm{Z}_i\hat{\bm{u}}_i \\
%			  &=  \bm{X}_{i,*}\hat{\bm{\beta}} + \bm{Z}_i\bm{G Z_i'V_i^{-1}(y_i-X_{i}\hat{\bm{\beta}})} \\
%			  &=  \bm{X}_{i,*}\hat{\bm{\beta}} + (\bm{Z_iG Z_i'+\Sigma_i})\bm{V_i^{-1}(y_i-X_{i}\hat{\bm{\beta}})} - \bm{\Sigma_i}\bm{V_i^{-1}(y_i-X_{i}\hat{\bm{\beta}} )}\\
%&= \bm{(\Sigma_i V_i^{-1}) \cdot X_{i,*}\hat{\bm{\beta}}} + (\bm{I_{T_i} - \Sigma_i V_i^{-1}}) \cdot \bm{y}_i
%\end{aligned}
%$$
%\end{theorem}

\begin{frame}[allowframebreaks]
\frametitle{Interval Estimation}

\begin{theorem}[Estimation of $Var(\hat{\bm{\beta}}^{OLS})$]
If $X$ is of full rank and assume that $\bm{\epsilon} \sim \mathcal{N}(\bm{0, \Sigma}=\sigma^2 I_m)$, then $Var(\hat{\bm{\beta}})$ is given as follows:
$$
Var(\hat{\bm{\beta}}^{OLS})=\sigma^2 \bm{(X'X)}^{-1} .
$$
\end{theorem}

\begin{proof} 
$\hat{\bm{\beta}}^{OLS}=\bm{(X'X)^{-1} X'y} = \bm{(X'X)^{-1} X'(X\beta + \epsilon)} = \bm{\beta + (X'X)^{-1} X'\epsilon}$, so we have
$$
\begin{aligned}
Var(\hat{\bm{\beta}}^{OLS}) &= \bm{(X'X)^{-1}X'}  Var( \bm{ \epsilon}) \bm{X(X'X)^{-1} } \\
&= \bm{(X'X)^{-1}X'}   \sigma^2 I_m\bm{X(X'X)^{-1} } = \sigma^2 \bm{(X'X)}^{-1}.
\end{aligned}
$$

\end{proof}

\framebreak

\begin{itemize}
\item Note that $Var(\hat{\bm{\beta}})$ is a covariance matrix so that the diagonal components are $\left(Var(\hat{\beta_1}), Var(\hat{\beta_2}), \ldots, Var(\hat{\beta_p}) \right)$,
where
$$
Var(\hat{\beta_j}) = \left[Var(\hat{\bm{\beta}})\right]_{j,j} = \sigma^2 \left[\bm{(X'X)}^{-1}\right]_{j,j} ,
$$
and
$$
\widehat{Var}(\hat{\beta_j})  =\widehat{\left[Var(\hat{\bm{\beta}})\right]}_{j,j} =\hat{\sigma}^2 \left[\bm{(X'X)}^{-1}\right]_{j,j}.
$$

\end{itemize}

\framebreak
We are interested whether $j^{th}$ covariate has significant impact on $\mathbf{y}$, which can be expressed as the following hypotheses:
$$
H_0: \beta_j=0 \qquad \text{VS} \qquad H_1: \beta_j \ne 0.
$$
Firstly, we define $t_j = {\hat{\beta}_j}/{\sqrt{\widehat{Var}(\hat{\beta_j})}}$.
We reject $H_0$ (equivalently consider $\beta_j$ as significant) at $95\%$ level if 
\begin{itemize}
	\item $\mathbb{P} \left(t_{n-p;0.975}  \geq \left|t_j \right| \right) \leq 0.05$, or
	\item $\left|t_j\right| \geq t_{n-p;0.975} \overset{n \rightarrow \infty}\simeq z_{0.975} = 1.96$,  or
	\item $0 \not\in (\hat{\beta}_j - t_{n-p;0.975} \cdot \sqrt{\widehat{Var}(\hat{\beta_j})}, \ \hat{\beta}_j + t_{n-p;0.975}\cdot \sqrt{\widehat{Var}(\hat{\beta_j})}  )$.

\end{itemize}

Note that $t_j \sim t(n-p)$ under $H_0$.

\end{frame}



\begin{frame}{Residual Sum of Squares (RSS)}

The residual sum of squares is:
\[
\mathrm{RSS} = \sum_{i=1}^n (y_i - \hat{y}_i)^2 = (y - \hat{y})^\top (y - \hat{y}) = \hat{\sigma}^2 (n-p)
\]

Using $\hat{y} = H y$ where $H = X'(X'X)^{-1}X$:
\[
\mathrm{RSS} = (y - Hy)^\top (y - Hy)
= y^\top (I - H)^\top (I - H) y
\]

Since $H$ is symmetric and idempotent ($H = H^\top = H^2$):
\[
\boxed{
\mathrm{RSS} = y^\top (I - H)  y = \hat{\sigma}^2 (n-p)
}
\]

It is also known that $\frac{\mathrm{RSS}}{\sigma^2} \sim \chi^2_{n-p}
$ as $\mathrm{Rank}(I-H)=n-p$.

\end{frame}


\begin{frame}{Explained  Sum of Squares (ESS)}

The explained sum of squares is:
\[
\mathrm{ESS} = \sum_{i=1}^n (\hat{y}_i - \bar{y})^2 = (\hat{y} - \bar{y}\mathbf{1})^\top (\hat{y} - \bar{y}\mathbf{1})
= \hat{y}^\top (I - J)\hat{y},
\]

where $\mathbf{1}$ be the vector of ones and $J = \frac{1}{n}\mathbf{1}\mathbf{1}^\top$.


Substitute $\hat{y} = Hy$:
\[
\mathrm{ESS} = y^\top H (I - J) H y
\]

If the model includes an intercept (so that $HJ = J$), this simplifies to:
\[
\boxed{
\mathrm{ESS} = y^\top (H - J) y
}
\]

It is also known that $\frac{\mathrm{ESS}}{\sigma^2} \sim \chi^2_{p-1}
$ {\color{red}provided $H_0: \beta_2=\cdots=\beta_p=0$ is true}.


\end{frame}


\begin{frame}{Significance of the regression model}

We wonder if the fitted regression model is (statistically) meaningful, which is equivalent to test the following:
$$
H_0: \beta_2=\cdots=\beta_p=0 \quad \text{ vs } \quad  H_1: H_0 \text{ is not true}
$$
\begin{itemize}
    \item Intuitively, we expect that if $\mathrm{ESS}$ is larger (i.e., simple $\bar{y}$ is not enough to explain the variations in $y$), then $H_0$ is less likely to be true. 
    \item So we want to reject $H_0$ if $\mathrm{ESS}$ is `large enough' compared to $\mathrm{RSS}$.

Wait..... how large is `large enough'?

\end{itemize}

\end{frame}


\begin{frame}{Significance of the regression model}

We wonder if the fitted regression model is (statistically) meaningful, which is equivalent to test the following:
$$
H_0: \beta_2=\cdots=\beta_p=0 \quad \text{ vs } \quad  H_1: H_0 \text{ is not true}
$$
Note that we know that $\mathrm{RSS}$ and $\mathrm{ESS}$ are independent (why?),  and $\frac{\mathrm{RSS}}{\sigma^2} \sim \chi^2_{n-p}$ and $\frac{\mathrm{ESS}}{\sigma^2} \sim \chi^2_{p-1}
$ {\color{red}provided $H_0$ is true}.

Therefore, under $H_0$, we know that
$$
{F}_{all}=\frac{\mathrm{ESS} / (p-1)}{\mathrm{RSS} / (n-p)} \sim F_{p-1, n-p},
$$
and reject $H_0$ at the significance level $\alpha$ if
$$
{F}_{all} > F_{p-1, n-p; \alpha}.
$$

\end{frame}


\begin{frame}{Significance of the regression coefficient}

We consider if the $j^{th}$ covariate is (statistically) meaningful, which is equivalent to test the following:
$$
H_0: \beta_j=0 \quad \text{ vs } \quad  H_1: H_0 \text{ is not true}
$$
Similarly, it is known that 
$ \color{blue}
(\hat{\beta}_j -\beta_j)/ \sqrt{Var(\hat{\beta}_j)}={\hat{\beta}_j}/{\sigma \sqrt{\left[\bm{(X'X)}^{-1}\right]_{j,j}}} \sim \mathcal{N}(0,1)$ under $H_0$ so that $(\hat{\beta}_j / Var(\hat{\beta}_j) )^2 \sim \chi^2_1$, 
 and $\hat{\beta}_j$ and $\color{red} \mathrm{RSS}$ are independent (why?) provided $H_0$ is true.

Therefore, under $H_0$, we know that
$$
t^2_j = \frac{\color{blue} {\hat{\beta}^2_j}/{\sigma^2 {\left[\bm{(X'X)}^{-1}\right]_{j,j}}}}{\color{red} \hat{\sigma}^2 /\sigma^2 }={F}_j=\frac{\color{blue}(\hat{\beta}_j / Var(\hat{\beta}_j) )^2}{\color{red}\mathrm{RSS} / (\sigma^2(n-p))}  \sim F_{1, n-p},
$$
and reject $H_0$ at the significance level $\alpha$ if
$$
t^2_j={F}_j > F_{1, n-p; \alpha}.
$$

\end{frame}

\begin{frame}{Total Sum of Squares (TSS)}
By definition, \[
\mathrm{TSS} = \sum_{i=1}^n (y_i - \bar{y})^2
= y^\top (I - J) y
\] and
\[
\boxed{
\mathrm{TSS} = \mathrm{ESS} + \mathrm{RSS}
}
\]



or equivalently:
\[
y^\top (I - J) y = y^\top (H - J) y + y^\top (I - H) y
\]

\end{frame}

\begin{frame}{Coefficient of determination: $R^2$}
As mentioned before, if $\mathrm{ESS}$ is larger, it implies that the regression model will be more meaningful. So we define the following measure accordingly:
$$
R^2 = \frac{\mathrm{ESS}}{\mathrm{TSS}} =1- \frac{\mathrm{RSS}}{\mathrm{TSS}}
$$
Since $1-R^2 = \frac{\mathrm{RSS}}{\mathrm{TSS}}$, we also have
$$
\hat{F}=\frac{\mathrm{ESS} / (p-1)}{\mathrm{RSS} / (n-p)} = \frac{R^2 / (p-1)}{(1-R^2) / (n-p)} = \frac{n-p}{p-1} \cdot \frac{R^2}{1-R^2}.
$$

\end{frame}


\end{document}